\subsection{Dedekind 整环上的模}

我们现在假设 A 是 Dedekind 整环；于是我们知道理想 \(p \in P\) 都是极大的，且它们是 A 的仅有的 \(\neq 0\) 素理想（ \(\S 2\) ，第 1 节）；群 D(A) 等同于 A 的 \(\neq 0\) 分式理想所成的群 I(A)。

命题 21. 设 A 是 Dedekind 整环。每个伪零 A-模都是零。每个伪单射（相应地，伪满射、伪双射、伪零）A-模同态都是单射（相应地，满射、双射、零）。

第一个断言已证明（第 4 节，例 1）；其余的由它立即得出。

命题 22. 设 A 是 Dedekind 整环，M 是有限生成 A-模。下列性质等价：

\begin{enumerate}
\def\labelenumi{(\alph{enumi})}
\item
  M 无扭。
\item
  M 自反。
\item
  M 是投射的。
\end{enumerate}

我们已经知道（不对整环 A 作任何假设）：(b) 蕴含 (a)（第 2 节，注 1），(c) 蕴含 (b)（《代数》第 II 章，§ 2，第 7 节，命题 14 的推论 4）。若 M 无扭，它就等同于 \(\mathbf{V} = \mathbf{M} \otimes_{\mathbf{A}} \mathbf{K}\) 关于 A 的一个格；于是 \(\mathbf{M}_{\mathfrak{p}}\) 是自由 \(\mathbf{A}_{\mathfrak{p}}\) -模，对每个极大理想 \(\mathfrak{p} \in \mathbb{P}\) 皆然，因为 \(\mathbf{A}_{\mathfrak{p}}\) 是主理想整环。于是结论由第 II 章，§ 5，第 2 节，定理 1 (b) 得出。

推论. 设 M 是有限生成 A-模，T 是它的扭子模。则 T 是 M 的直因子。

由于 M/T 无扭且有限生成，由命题 22 它是投射的，于是推论由《代数》第 II 章，§ 2，第 2 节，命题 4 得出。

命题 23. 设 A 是 Dedekind 整环，T 是有限生成扭 A-模。存在两个有限族 \((n_{i})_{i\in I}\) 与 \((\mathfrak{p}_i)_{i\in I}\) ，其中诸 \(n_i\) 是 \(\geqslant 1\) 的整数，诸 \(\mathfrak{p}_i\) 是 P 的元素，使得 T 同构于直和 \(\bigoplus_{i\in I} (\mathrm{A} / \mathfrak{p}_{i}^{n_i})\) 。此外，族 \((n_{i})_{i\in I}\) 与 \((\mathfrak{p}_i)_{i\in I}\) 在相差指标集的一个双射的意义下是唯一的。

这由第 4 节的定理 5 得出，并注意到此处伪同构就是同构。

命题 24. 设 A 是 Dedekind 整环，M 是秩 \(n \geqslant 1\) 的有限生成无扭 A-模。则存在 A 的理想 \(d \neq 0\) ，使 M 同构于模 \(A^{n-1}\) 与 d 的直和。此外，理想 d 的类由这一条件唯一确定。

第 9 节的定理 6 表明存在 M 的自由子模 L，使 M/L 同构于 A 的理想 a。若 a = 0，我们取 d = A。否则，a 的秩为 1，于是 \(L = A^{n-1}\) 且 a 是投射模（命题 22）；因此 M 同构于 L 与 a 的直和（《代数》第 II 章，§2，第 2 节，命题 4），这就证明了命题的第一部分。此外，由第 7 节命题 16 (i)、(iv) 与 (v) 得 \(c(\mathbf{M}) = c(\mathfrak{d})\) ，由此得 d 的类的唯一性。

\textbf{注}\par

\begin{enumerate}
\def\labelenumi{(\arabic{enumi})}
\item
  命题 23、命题 24 以及命题 22 的推论完全确定了有限生成 A-模的结构。命题 24 表明，无扭的有限生成 A-模在同构意义下由它的秩及附加于它的除子类确定。
\item
  可以证明：在 Dedekind 整环上，非有限生成的投射模必然是自由的（习题 21），且投射模的每个子模都是投射的（习题 20）。
\end{enumerate}


